\begin{textsample}{POS Dim 3 – persona\_grok – Score 18.00 – t845\_grok.txt}  \label{ex:f3_pos_034}
In Brazil, a major scandal has exposed deep flaws in the \textbf{meat} \textbf{industry}.
Police investigations revealed that big \textbf{companies} like JBS and BRF were bribing inspectors to certify \textbf{meat} that wasn't safe for \textbf{consumption}.
Practices included \textbf{using} \textbf{chemicals} to hide the smell and appearance of rotten \textbf{meat}, mixing pigs ' heads into sausages, and even adding cardboard to chicken \textbf{products}.
The fallout was swift : countries including China and the European Union, along with others, imposed bans on Brazilian \textbf{meat} imports to protect \textbf{public} \textbf{health}.
This isn't just a one-off issue—it 's a symptom of broader \textbf{problems} in global industrial \textbf{meat} \textbf{production}.
The \textbf{system} is tied to outbreaks of diseases and significant environmental damage, such as widespread deforestation.
Industrial \textbf{meat} \textbf{production} contributes around 14 % of global \textbf{greenhouse} \textbf{gas} \textbf{emissions}.
At Greenpeace, we 're calling for change. \textbf{Shifting} to plant-rich \textbf{diets} and eating less \textbf{meat} overall can make a real difference.
We also \textbf{need} to hold governments accountable for regulating the \textbf{industry} better.
Our campaign pushes for \textbf{meat} that 's sourced locally and sustainably, ensuring it 's produced in ways that respect the planet.
In a recent update to this story, it 's come to light that these \textbf{companies} have been purchasing \textbf{cattle} raised on lands that were illegally deforested, further fueling environmental destruction.

% matched lemmas: cattle, chemical, company, consumption, diet, emission, gas, greenhouse, health, industry, meat, need, problem, product, production, public, shift, system, use
\end{textsample}
